Although nicely generic water quality constraints are presented above, i.e. \eqref{eq:inlet_treatability} to \eqref{eq:outlet_water_quality},
we shall devote this section to interpret the quality parameters $p \in \mathcal P$ and how they are treated in plants.
We confine ourselves to ``removal process'', though our model can be extended to include more treatment processes, e.g. ``chemical dosing'', and thus, this section will be adapted to be consistent with the existing parameter notation while preserving the linear structure of the MILP.

\subsection{Inlet treatability and Outlet water quality}

In \eqref{eq:inlet_treatability}, we propose acceptance criteria with respect to the incoming $p \in \mathcal P$ at every plant $t \in \mathcal T$.
Due to linearity assumption, we insist these quality parameters can be blended approximately linearly by volume, and the source value can be used directly from $Q_{sp}$.
In this constraint, we avoid division by decision variables $b_{st}$, which result in a non–linear program, we write the constraint in load form.
It can be understood as a volume–weighted average quality, rather than pure quality constraints.

In \eqref{eq:outlet_water_quality}, we restrict the outgoing $p \in \mathcal P$ after undergoing the treatment processes at $t \in \mathcal T$ to be within a certain regulatory range.
Likewise, to avoid division by decision variables $b_{st}$, we express the constraints in load form, or equivalently, volume–weighted average.

\subsection{Turbidity}

Turbidity measures the cloudiness of water and is associated with suspended particles such as sediment, organic matter, algae, clay, or runoff material.
In water supply modelling, turbidity is important because high turbidity may increase treatment requirements and affect disinfection performance \cite{WHOTurbidity2011, Matos2024Turbidity, UMNTurbidity}.

For turbidity, the quality parameter is represented by setting $p=\mathrm{turb}$.
Thus, $Q_{s,\mathrm{turb}}$ denotes the turbidity value of source $s$, while $\underline{Q}_{\mathrm{turb}}$ and $\overline{Q}_{\mathrm{turb}}$ denote the lower and upper turbidity limits.

Conveniently, the turbidity requirement is handled directly by replacing the sub–index $p$ with $\mathrm{turb}$.
For example, \eqref{eq:inlet_treatability} becomes:

\begin{equation*}
    \underline{I}_{t,\mathrm{turb}} ~\sum_{s \in \mathcal{S}} b_{st}
    \le \sum_{s \in \mathcal{S}}
    Q_{s,\mathrm{turb}} ~b_{st} \le \overline{I}_{t,\mathrm{turb}} ~\sum_{s \in \mathcal{S}} b_{st}, \quad \forall t \in \mathcal{T}.
\end{equation*}

\subsection{Alkalinity}

Alkalinity measures the ability of water to neutralise acid \cite{ADWG2011, TestNeeds2025Alkalinity, PerfectWater2025Alkalinity}.
It represents the buffering capacity of water and is relevant to $\mathrm{pH}$ stability, corrosion control, treatment performance, remineralisation, and water stability in distribution systems \cite{ADWG2011, PerfectWater2025Alkalinity}.

For alkalinity, the quality parameter is represented by setting $p = \mathrm{alk}$.
Thus, $Q_{s,\mathrm{alk}}$ denotes the alkalinity value of source $s$, while $\underline{Q}_{\mathrm{alk}}$ and $\overline{Q}_{\mathrm{alk}}$ denote the lower and upper alkalinity limits.

Conveniently, the alkalinity requirement is handled directly by replacing the sub–index $p$ with $\mathrm{alk}$.
For example, \eqref{eq:inlet_treatability} becomes:

\begin{equation*}
    \underline{I}_{t,\mathrm{alk}} ~\sum_{s \in \mathcal{S}} b_{st}
    \le \sum_{s \in \mathcal{S}}
    Q_{s,\mathrm{alk}} ~b_{st} \le \overline{I}_{t,\mathrm{alk}} ~\sum_{s \in \mathcal{S}} b_{st}, \quad \forall t \in \mathcal{T}.
\end{equation*}

We highlight here that, according to \cite{who2004water}, removal process can directly reduce alkalinity, though it can be increased by an adjustment process.

\subsection{pH transformation}

Unlike Turbidity and Alkalinity, pH cannot be treated in any removal processes, but rather in other treatment processes, e.g. chemical dosing,  coagulation and lime softening \cite{who2004water}.

More importantly, $\mathrm{pH}$ values are not averaged directly because direct $\mathrm{pH}$ averaging can produce a scientifically misleading blended value.
Therefore, in our mathematical model, we consider its corresponding hydrogen ion.

\begin{equation}\label{eq:ph_definition}
    \mathrm{pH} = -\log_{10}([\mathrm{H}^{+}]),
\end{equation}
where $[\mathrm{H}^{+}]$ represents hydrogen ion concentration \cite{EPA2026pH}.

To keep the MILP linear, $\mathrm{pH}$ is handled as a preprocessing transformation \cite{WHO2007pH}.
For the $\mathrm{pH}$ quality parameter, $Q_{s,\mathrm{pH}}$ stores the transformed hydrogen ion concentration obtained from the raw source $\mathrm{pH}$ value:
\begin{equation}\label{eq:ph_source_transformation}
    Q_{s,\mathrm{pH}} = 10^{-\mathrm{pH}_s}.
\end{equation}

The symbols $\underline{\mathrm{pH}}$ and $\overline{\mathrm{pH}}$ represent the raw regulatory lower and upper limits in $\mathrm{pH}$ units.
These raw limits are converted once during data preparation before being passed into the MILP.
If the acceptable raw $\mathrm{pH}$ range is:
\begin{equation}\label{eq:raw_ph_range}
    \underline{\mathrm{pH}} \leq
    \mathrm{pH}
    \leq
    \overline{\mathrm{pH}},
\end{equation}
then the corresponding transformed bounds stored for the $\mathrm{pH}$ quality parameter are:
\begin{equation}\label{eq:ph_bound_transformation}
    \underline{Q}_{\mathrm{pH}} = 10^{-\overline{\mathrm{pH}}},
    \qquad
    \overline{Q}_{\mathrm{pH}} = 10^{-\underline{\mathrm{pH}}}.
\end{equation}

The direction of the bounds changes because $\mathrm{pH}$ and hydrogen ion concentration are inversely related \cite{Fondriest2025pH}.
A higher $\mathrm{pH}$ corresponds to a lower hydrogen ion concentration, while a lower $\mathrm{pH}$ corresponds to a higher hydrogen ion concentration.

As mentioned earlier, pH cannot be handled by removal process, thereby, inapplicable to constraints from \eqref{eq:removal_process_initial} to \eqref{eq:outlet_water_quality}.
Nonetheless, as an incoming quality parameter at $t \in T$, the constraint \eqref{eq:inlet_treatability} is still valid. Formally, this is:

\begin{equation*}
    \underline{I}_{t,\mathrm{pH}} ~\sum_{s \in \mathcal{S}} b_{st}
    \le \sum_{s \in \mathcal{S}}
    Q_{s,\mathrm{pH}} ~b_{st} \le \overline{I}_{t,\mathrm{pH}} ~\sum_{s \in \mathcal{S}} b_{st}, \quad \forall t \in \mathcal{T}.
\end{equation*}

\subsection{Blended pH recovery after optimisation}

The optimisation model constrains $\mathrm{pH}$ through the transformed value $Q_{s,\mathrm{pH}}$.
After the solver finds the optimal source to plant flows, the blended hydrogen ion concentration at plant $t$ is calculated as:
\begin{equation}\label{eq:blended_ph_transformed}
    Q^{\mathrm{blend}}_{t,\mathrm{pH}} = \frac{ \sum_{s \in \mathcal{S}} Q_{s,\mathrm{pH}} b_{st} }{ \sum_{s \in \mathcal{S}} b_{st} }.
\end{equation}

The blended $\mathrm{pH}$ is then recovered in post-processing as:
\begin{equation}\label{eq:blended_ph_recovery}
    \mathrm{pH}^{\mathrm{blend}}_{t}
    =
    -\log_{10}
    \left(
    Q^{\mathrm{blend}}_{t,\mathrm{pH}}
    \right).
\end{equation}

This recovered $\mathrm{pH}$ value is not an MILP decision variable.
It is a post-solve reporting value used for validation, interpretation, dashboard display, and output generation.

\textcolor{red}{\textbf{TODO:} update these to reflect new removal formulation i.e. last assumption is now incorrect and there should be more assumptions about the removal processes.}

\textcolor{blue}{I think you meant the Section 6.4? Because I think the Section 6.5 is sound, and I modified few things in Section 6.4, also other sections before 6.4. One more thing, I added a separate subsection for 6.7 Assumptions for removal process.}

\subsection{Linearity assumptions for water quality}

The water quality formulation relies on simplifying assumptions that preserve linearity in the MILP.
The key assumptions are:
\begin{enumerate}[label=(\roman*), align=left, widest=viii]
    \item Water from selected sources is completely mixed before quality is assessed.
    \item The model represents a single time period, so temporal variation in water quality is not modelled.
    \item Source quality values are fixed input parameters for the representative period.
    \item Alkalinity and turbidity are approximated as volume-weighted linear blends.
    \item Raw $\mathrm{pH}$ values are not averaged directly; $\mathrm{pH}$ is transformed before optimisation.
    \item No nonlinear chemical reactions between blended sources are modelled.
    \item Water quality is checked at each treatment plant, using the blend of its incoming sources.
\end{enumerate}

These assumptions allow the water quality constraints to remain linear and suitable for inclusion in the MILP formulation.

\subsection{Assumptions for removal processes}
In practices, several designs for removal processes exist, see \cite{lower_murray_water2023, zouboulis2005improvement}. Therefore, we shall propose the following assumptions for convenience and simplicity.
\begin{itemize}
    \item Removal process $\mathcal R_t$ is conducted in a sequence.
    \item Maximum removal efficiency $H_{trp}$ is known. \textcolor{red}{I'm thinking about a weaker assumption, not maximum but a fixed efficiency is known?}
    \item Uniform treatment: $\sum_{s \in \mathcal S} b_{st}$ receives the same treatment.
    \item No water loss after passing through $r_i \in \mathcal R_t$.
\end{itemize}

